\documentclass[a4paper,12pt]{article}
\usepackage[T2A]{fontenc}
\usepackage[utf8]{inputenc}
\usepackage[russian]{babel}
\usepackage{amsmath,amssymb,mathtools}
\usepackage{DejaVuSansCondensed}
\renewcommand{\familydefault}{\sfdefault}
\usepackage{icomma}
\usepackage{geometry}
\geometry{left=20mm,right=20mm,top=21mm,bottom=22mm}
\usepackage{microtype}
\usepackage{tikz}
\usetikzlibrary{calc,arrows.meta,positioning,shapes.geometric}
\usepackage{pgfplots}
\pgfplotsset{compat=1.18}
\usepackage{tabularx,booktabs}
\usepackage{enumitem}
\usepackage{fancyhdr}
\usepackage{setspace}
\usepackage{xcolor}
\usepackage{array}
\usepackage{hyperref}
\hypersetup{hidelinks,pdftitle={Зависимости величин, формулы и графики},pdfauthor={Артём Александрович Лёвин}}
\definecolor{accent}{HTML}{087B7B}
\definecolor{darkaccent}{HTML}{075353}
\definecolor{textgray}{HTML}{343B40}
\definecolor{muted}{HTML}{6E7B7B}
\definecolor{pale}{HTML}{E7F1F0}
\color{textgray}
\onehalfspacing
\setlength{\parindent}{0pt}
\setlength{\parskip}{6pt}
\setlist[itemize]{leftmargin=1.5em,itemsep=2pt,topsep=3pt}
\setlist[enumerate]{leftmargin=1.6em,itemsep=5pt,topsep=4pt}
\renewcommand{\arraystretch}{1.28}
\newcommand{\term}[1]{\textcolor{darkaccent}{\textbf{#1}}}
\newcommand{\markcheck}{\(\square\)\,}
\newcommand{\lessonsection}[1]{\vspace{5pt}{\large\bfseries\color{darkaccent}#1}\par\vspace{3pt}}
\pagestyle{fancy}
\fancyhf{}
\fancyhead[L]{\small\color{muted}Математика \(\cdot\) 7 класс}
\fancyhead[R]{\small\color{muted}21.10.26}
\fancyfoot[C]{\small\color{muted}\thepage}
\renewcommand{\headrulewidth}{0.25pt}
\renewcommand{\footrulewidth}{0pt}
\setlength{\headheight}{16pt}
\begin{document}
\thispagestyle{empty}
\begin{center}
\vspace*{22mm}
{\large\color{muted}\textbf{МАТЕМАТИКА · 7 КЛАСС}}\par
\vspace{24mm}
{\fontsize{34}{39}\selectfont\bfseries\color{darkaccent}ЧЕК-ЛИСТ}\par
\vspace{10mm}
{\LARGE\bfseries Зависимости величин,\par формулы и графики}\par
\vspace{13mm}
{\large От текста задачи к формуле,\par таблице значений и точкам на графике}\par
\vfill
{\large\textbf{Кирилл Зиновьев}}\par
\vspace{3mm}
{\normalsize 21 октября 2026 года}\par
\vspace{17mm}
{\small\color{muted}Лёвин Артём Александрович\par эксклюзивно для Кирилла Зиновьева}\par
\vspace{10mm}
\begin{tikzpicture}
\draw[accent,line width=1.25pt] (0,0) .. controls (3,0.55) and (8,-0.65) .. (13,0.15);
\end{tikzpicture}
\end{center}
\clearpage

\lessonsection{1. Как перевести условие задачи на язык математики}
Главная формулировка занятия: \term{«зависимость величины \(A\) от величины \(B\)»}. Это значит, что \(A\) вычисляется по выбранному значению \(B\):
\[
\boxed{\ A=f(B)\ }
\]
Буква \(B\) в скобках показывает, \term{от какой величины} зависит результат. Запись \(f\) означает правило вычисления. Например, \(s(t)\) читают как «путь как функция времени».

\begin{center}
\begin{tabularx}{\linewidth}{@{}>{\raggedright\arraybackslash}X>{\raggedright\arraybackslash}X@{}}
\toprule
\textbf{В условии написано} & \textbf{Как расположить переменные}\\
\midrule
Зависимость пути \(s\) от времени \(t\) & \(s\) слева, \(t\) справа: \(s(t)=\ldots\)\\
Зависимость времени \(t\) от пути \(s\) & \(t\) слева, \(s\) справа: \(t(s)=\ldots\)\\
Зависимость массы \(m\) от объёма \(V\) & \(m\) слева, \(V\) справа: \(m(V)=\ldots\)\\
\bottomrule
\end{tabularx}
\end{center}

\markcheck Я умею определить зависимую и независимую величины по словам «зависимость \(\ldots\) от \(\ldots\)».

\lessonsection{2. Путь, скорость и время}
При равномерном движении справедлива формула
\[
\boxed{s=vt},\qquad
\begin{aligned}
s&\text{ --- путь,}\quad
v\text{ --- скорость,}\quad
t\text{ --- время.}
\end{aligned}
\]
Если автомобиль едет со скоростью \(80\) км/ч, то зависимость пути от времени имеет вид
\[
\boxed{s(t)=80t},\qquad t\ge 0.
\]
Время выбираем самостоятельно, а путь \term{рассчитываем} по формуле. При \(t=0\) путь равен нулю; при \(t=1\) ч --- \(80\) км; при \(t=2\) ч --- \(160\) км.

\begin{center}
\begin{tabular}{@{}lcccc@{}}
\toprule
\(t\), ч & \(0\) & \(1\) & \(2\) & \(3\)\\
\midrule
\(s\), км & \(0\) & \(80\) & \(160\) & \(240\)\\
\bottomrule
\end{tabular}
\end{center}
\vspace{-2pt}
\begin{center}
\begin{tikzpicture}
\begin{axis}[
 width=10.2cm,height=3.75cm,
 xmin=-0.15,xmax=3.3,ymin=-10,ymax=255,
 axis lines=middle,axis line style={->},
 xlabel={$t$, ч},ylabel={$s$, км},
 xlabel style={at={(axis description cs:1.02,0.08)},anchor=west,font=\small},
 ylabel style={at={(axis description cs:0.1,0.98)},anchor=north west,font=\small},
 xtick={0,1,2,3},ytick={0,80,160,240},
 grid=major,major grid style={draw=gray!20},
 tick label style={font=\scriptsize},clip=false]
 \addplot[darkaccent,very thick,domain=0:3,samples=2] {80*x};
 \addplot[only marks,mark=*,mark size=2pt,color=accent] coordinates {(0,0) (1,80) (2,160) (3,240)};
\end{axis}
\end{tikzpicture}
\end{center}
\vspace{-4pt}


\clearpage
\lessonsection{3. Почему слева бывает не только \texorpdfstring{\(s\)}{s}}
Если время движения неизменно и равно \(6\) ч, а скорость изменяется, то
\[
\boxed{s(v)=6v}.
\]
Например, \(s(65)=6\cdot65=390\) км. Если за \(6\) ч проехали \(363\) км, то \(6v=363\), откуда \(v=60{,}5\) км/ч.

Если просят, наоборот, \term{зависимость времени от пути} при скорости \(12\) км/ч, то из \(s=vt\) выражаем время:
\[
\boxed{t(s)=\frac{s}{12}},\qquad s\ge0.
\]
\markcheck Я записываю слева именно ту величину, зависимость которой требуется выразить.

\lessonsection{4. Масса и объём: та же идея в другой задаче}
Масса однородного материала вычисляется по формуле \(m=\rho V\), где \(\rho\) --- плотность, \(V\) --- объём. Для пробки плотностью \(0{,}18\) г/см\(^3\) получаем
\[
\boxed{m(V)=0{,}18V},\qquad V\ge0.
\]
\textbf{Известен объём:} \(V=240\) см\(^3\) \(\Rightarrow\)
\(m=0{,}18\cdot240=43{,}2\) г.

\textbf{Известна масса:} \(m=64{,}8\) г \(\Rightarrow\)
\(0{,}18V=64{,}8\), следовательно, \(V=360\) см\(^3\).

\markcheck Я понимаю, когда подставлять число в формулу, а когда сначала составлять уравнение.

\lessonsection{5. Как строить график по формуле}
\begin{enumerate}[label=\textcolor{accent}{\arabic*.},leftmargin=1.45em]
\item Запишите формулу зависимости и допустимые значения переменной.
\item Выберите значения \term{независимой} переменной и вычислите соответствующие результаты.
\item Начертите и подпишите оси: \(x\) (или \(t\)) --- \term{горизонтальная}, \(y\) (или \(s\)) --- \term{вертикальная}.
\item Укажите масштаб \term{отдельно для каждой оси}; нанесите точки с координатами \((x;y)\).
\item Соедините точки с учётом вида зависимости: для \(s=80t\) получается луч, для \(y=x(x-3)\) --- плавная кривая.
\end{enumerate}
\markcheck Я отмечаю точку \((-1;4)\): сначала на \(1\) влево по \(x\), затем на \(4\) вверх по \(y\).

\clearpage
\lessonsection{6. Таблица и график функции \texorpdfstring{\(y=x(x-3)\)}{y=x(x-3)}}
Пусть \(x\) изменяется от \(-2\) до \(2\), включая границы:
\[
\boxed{y=x(x-3)},\qquad -2\le x\le2.
\]
\textbf{Пример подстановки:} при \(x=-1{,}5\) получаем
\[
y(-1{,}5)=(-1{,}5)(-1{,}5-3)=(-1{,}5)(-4{,}5)=6{,}75.
\]
\vspace{-3pt}
\begin{center}
\small
\begin{tabular}{@{}lrrrrrrrrr@{}}
\toprule
\(x\) & \(-2\) & \(-1{,}5\) & \(-1\) & \(-0{,}5\) & \(0\) & \(0{,}5\) & \(1\) & \(1{,}5\) & \(2\)\\
\midrule
\(y\) & \(10\) & \(6{,}75\) & \(4\) & \(1{,}75\) & \(0\) & \(-1{,}25\) & \(-2\) & \(-2{,}25\) & \(-2\)\\
\bottomrule
\end{tabular}
\end{center}

\begin{center}
\begin{tikzpicture}
\begin{axis}[
 width=12.0cm,height=8.2cm,
 xmin=-2.6,xmax=2.55,ymin=-3.25,ymax=11.2,
 axis lines=middle,axis line style={->,line width=0.75pt},
 xlabel={$x$},ylabel={$y$},
 xlabel style={at={(axis description cs:0.99,0.22)},anchor=north east},
 ylabel style={at={(axis description cs:0.54,0.99)},anchor=north east},
 xtick={-2,-1,0,1,2},ytick={-2,0,2,4,6,8,10},
 grid=both,minor tick num=1,
 major grid style={draw=gray!25},minor grid style={draw=gray!12},
 tick label style={font=\scriptsize},
 clip=false]
 \addplot[darkaccent,very thick,smooth,domain=-2:2,samples=201] {x*(x-3)};
 \addplot[only marks,mark=*,mark size=2pt,color=accent] coordinates {(-2,10) (-1.5,6.75) (-1,4) (-0.5,1.75) (0,0) (0.5,-1.25) (1,-2) (1.5,-2.25) (2,-2)};
 \node[anchor=south west,font=\small,text=darkaccent] at (axis cs:-1,4) {$(-1;4)$};
 \node[anchor=north west,font=\small,text=darkaccent] at (axis cs:1.5,-2.25) {$(1{,}5;-2{,}25)$};
\end{axis}
\end{tikzpicture}
\end{center}
\vspace{-5pt}
\markcheck Я умею вычислять значения функции при отрицательных и дробных \(x\), наносить точки и учитывать заданный промежуток.

\lessonsection{Проверьте себя: три частые ошибки}
\begin{itemize}
\item Записать \(s=12t\), когда просят зависимость \term{времени от пути}. Правильно: \(t(s)=s/12\).
\item Подставить \(x=-1{,}5\) лишь в одну из двух скобок. Подстановка должна быть полной.
\item Нарисовать график дальше заданного промежутка \([-2;2]\). За его пределами график в этой задаче не строят.
\end{itemize}

\clearpage
\begin{center}{\LARGE\bfseries\color{darkaccent}Алгоритм: от условия к формуле зависимости}\end{center}
\vspace{9mm}
\begin{center}
\tikzset{
 flow/.style={draw=accent,rounded corners=4pt,very thick,align=center,text width=10.4cm,minimum height=12mm,font=\normalsize,inner sep=5pt},
 flowend/.style={draw=darkaccent,rounded corners=12pt,line width=1.2pt,align=center,text width=9.2cm,minimum height=13mm,font=\normalsize,inner sep=5pt},
 decision/.style={diamond,draw=accent,very thick,aspect=2.7,align=center,text width=4.7cm,font=\small,inner xsep=2pt,inner ysep=4pt},
 smallflow/.style={draw=accent,rounded corners=3pt,line width=1.05pt,align=center,text width=4.3cm,minimum height=15mm,font=\small,inner sep=4pt},
 arrow/.style={-Latex,thick,color=darkaccent}
}
\begin{tikzpicture}[node distance=12mm and 16mm]
\node[flowend] (start) {Прочитайте:\quad «зависимость \(A\) от \(B\)»};
\node[flow,below=of start] (vars) {Определите: \(A\) --- результат, \(B\) --- независимая величина};
\node[flow,below=of vars] (law) {Найдите общую формулу связи величин\\(например, \(s=vt\) или \(m=\rho V\))};
\node[decision,below=15mm of law] (known) {Заданы числовые значения постоянных величин?};
\node[smallflow,below left=20mm and 11mm of known] (yes) {Подставьте эти числа вместо постоянных};
\node[smallflow,below right=20mm and 11mm of known] (no) {Сохраните неизвестные параметры буквами};
\node[flow,below=35mm of known] (write) {Выразите \(A\) через \(B\) и запишите \(A(B)=\ldots\)};
\node[flow,below=of write] (verify) {Проверьте единицы измерения и допустимые значения};
\node[flowend,below=of verify] (finish) {Формула готова: можно вычислять,\\составлять таблицу и строить график};
\draw[arrow] (start) -- (vars);
\draw[arrow] (vars) -- (law);
\draw[arrow] (law) -- (known);
\draw[arrow] (known.south west) -| node[pos=0.25,above left,font=\small]{да} (yes.north);
\draw[arrow] (known.south east) -| node[pos=0.25,above right,font=\small]{нет} (no.north);
\draw[arrow] (yes.south) |- ([yshift=6mm]write.north) -- (write.north);
\draw[arrow] (no.south) |- ([yshift=6mm]write.north);
\draw[arrow] (write) -- (verify);
\draw[arrow] (verify) -- (finish);
\end{tikzpicture}
\end{center}
\clearpage

\begin{center}
{\LARGE\bfseries\color{darkaccent}Тренировочный блок}\par
\vspace{2mm}{\small Путь А · 10 заданий · от простого к более сложному}\par
\end{center}
\textbf{Указание.} Записывайте формулы и необходимые вычисления. Графики выполняйте в тетради, подписывайте оси и масштаб.
\begin{enumerate}[label=\textcolor{darkaccent}{\textbf{\arabic*.}},leftmargin=1.5em,itemsep=11pt]
\item Автомобиль движется с постоянной скоростью \(20\) км/ч. Запишите формулу зависимости \textbf{пути от времени}. Укажите зависимую и независимую величины.
\item Для движения со скоростью \(80\) км/ч составьте таблицу значений \(s(t)\) при \(t=0;\ 0{,}5;\ 2;\ 3\) ч.
\item Автобус движется \(6\) ч с постоянной скоростью \(v\) км/ч. Запишите \(s(v)\). Найдите путь при \(v=58\) км/ч, а затем скорость, если путь равен \(420\) км.
\item Плотность пробки \(0{,}18\) г/см\(^3\). Запишите \(m(V)\), найдите массу при \(V=150\) см\(^3\) и объём при \(m=45\) г.
\item Велосипедист движется со скоростью \(12\) км/ч. Запишите формулу зависимости \textbf{времени от пути} \(t(s)\). Сколько часов потребуется для прохождения \(90\) км?
\end{enumerate}
\vfill
{\small\color{muted}Для всех задач на движение считайте, что скорость постоянна, а остановок нет.}
\clearpage

\begin{center}
{\LARGE\bfseries\color{darkaccent}Тренировочный блок · продолжение}\end{center}
\begin{enumerate}[label=\textcolor{darkaccent}{\textbf{\arabic*.}},leftmargin=1.5em,start=6,itemsep=10pt]
\item Дана функция \(s(t)=12t\), где \(0\le t\le3\). Составьте таблицу для \(t=0;\ 1;\ 2;\ 3\), постройте график. Определите по графику время, за которое пройдено \(30\) км.
\item Функция \(y=x(x-3)\) задана при \(-2\le x\le2\). Рассчитайте \(y\) для \(x=-2;\ -1;\ 0;\ 1;\ 2\). Принадлежит ли графику точка \((-1;4)\)?
\item Для той же функции найдите значения при \(x=-1{,}5;\ -0{,}5;\ 0{,}5;\ 1{,}5\). Покажите для одного значения подробную подстановку.
\item Используя результаты заданий 7 и 8, постройте график функции \(y=x(x-3)\) на отрезке \([-2;2]\). Укажите координаты самой нижней из нанесённых точек.
\item Ученик написал \(s=12t\) как ответ на требование «задайте формулой зависимость \textbf{времени от пути} при скорости \(12\) км/ч». Объясните ошибку, запишите верную формулу и найдите время для \(s=48\) км.
\end{enumerate}
\vspace{8mm}
\textbf{Перед сдачей проверьте:}
\begin{itemize}[label=\(\square\),leftmargin=1.4em,itemsep=4pt]
\item слева записана величина, зависимость которой требуется выразить;
\item в формулах подставлены все известные постоянные;
\item у вычислений есть нужные единицы измерения;
\item координаты записаны в порядке \((x;y)\), оси графиков подписаны;
\item на графике функции учтён заданный промежуток.
\end{itemize}
\clearpage

\begin{center}
{\LARGE\bfseries\color{darkaccent}Ответы для самопроверки}\par
\vspace{2mm}{\small Сначала выполните задания самостоятельно.}\par
\end{center}
\small
\begin{tabularx}{\linewidth}{@{}>{\bfseries\color{darkaccent}}p{0.9cm}X@{}}
\toprule
№ & Ответ\\
\midrule
1 & \(s(t)=20t\), \(t\ge0\). Независимая величина --- \(t\), зависимая --- \(s\).\\
2 & \(t:\ 0;\ 0{,}5;\ 2;\ 3\). \(s:\ 0;\ 40;\ 160;\ 240\) км.\\
3 & \(s(v)=6v\); \(348\) км; \(70\) км/ч.\\
4 & \(m(V)=0{,}18V\); \(27\) г; \(250\) см\(^3\).\\
5 & \(t(s)=\dfrac{s}{12}\); \(7{,}5\) ч.\\
6 & \(t:\ 0;\ 1;\ 2;\ 3\). \(s:\ 0;\ 12;\ 24;\ 36\) км. График --- отрезок прямой от \((0;0)\) до \((3;36)\); \(2{,}5\) ч.\\
7 & \(y:\ 10;\ 4;\ 0;\ -2;\ -2\). Да, точка \((-1;4)\) принадлежит графику.\\
8 & \(y:\ 6{,}75;\ 1{,}75;\ -1{,}25;\ -2{,}25\) (в указанном порядке).\\
9 & Плавная кривая через точки таблиц; самая нижняя из нанесённых точек: \((1{,}5;-2{,}25)\). График строится только при \(-2\le x\le2\).\\
10 & Требуется зависимость \(t\) от \(s\): \(t(s)=\dfrac{s}{12}\), \(t(48)=4\) ч.\\
\bottomrule
\end{tabularx}
\normalsize
\vfill
\begin{center}
\textcolor{darkaccent}{\rule{0.35\linewidth}{0.5pt}}\par
\small Вернитесь к пунктам чек-листа, которые вызвали затруднение.\par
\end{center}
\end{document}
